\section{讨论}
\label{sec:discussion}

\subsection{解读}

本研究的核心观察涉及两个解剖先验被引入同一细粒度分级骨干网络时的行为。
单独监督解剖部位,或单独监督弱病灶范围,均使复合分级风险劣于未经修改的
基线。同时提供二者则未继承该恶化,并取得四种配置中最佳观察的复合风险。这
是一种非可加的实证模式:联合配置的行为,不是任一单先验配置单独所能提示的。

我们可以给出一个合理的解释,但需说明本研究的实验并未分离其机制。每个先验
单独看,都增加了一个其最优解未必与细粒度分级对齐的训练目标。部位分类奖励
能区分扫查部位的特征,弱框注意奖励能集中于标注区域的特征;两个目标本身都
不直接关乎间隔纹理。孤立注入时,这样的信号可能挤占与分级相关的能力。而当
二者同时在场时,部位估计提供了病灶注意得以解释的解剖参照系,二者合起来对
证据的刻画可能比任何单独一个更完整。这一解释与观察一致,但并未被观察所
确立;要把多任务正则与预测注入分开,需要一个保留两个辅助损失、同时抑制
残差注入的配置。我们未运行该配置。

第二个观察更具直接可操作性。联合配置以 $5.9\%$ 的参数量与 $2.8\%$ 的算术
开销,产生了两个附加输出——六类解剖部位估计与弱病灶定位图——单张图像
中位延迟保持在 $4$~ms 以下。两个输出都仅由图像导出,推理时不需要任何元
数据。在筛查流程中,部位估计可用于自动核查标准化六部位协议是否被遵循,
定位图则为读片人提供了一个可与数值分数并列审视的对象。基线不具备这两种
能力中的任何一种,而且二者都不需要队列已有标注之外的额外标注。

分期分析反对用单一概括来描述困难所在。两个配置的平均绝对误差都在 F2 最大,
而召回率都在 F1 最低。因此这两个指标指认的最难分期并不相同。轻度纤维化的
超声表现可能比正常实质或晚期疾病更细微、更异质,这在临床上说得通,也与
观察到的 F1 召回率一致;我们将此作为一个合理的解读提出,而非已被证明的
解释。

最后,两种缩放约定的比较在复合终点上支持各向异性拉伸优于保持长宽比的信箱
式填充。我们据此将拉伸保留为所有建模配置的默认预处理,并在表~\ref{tab:main}
中报告信箱式对照,使该选择可见而非隐含。

\subsection{局限}

若干局限界定了本研究所能确立的范围。

\textbf{病因与机构范围。}队列全部由\emph{日本血吸虫}相关肝纤维化构成。
使解剖上下文在此显得合理的间隔状、网络状表型,不同于病毒性、酒精性与代谢
性肝病的弥漫性实质改变,我们的任何结果都不涉及那些病因。此外,虽然队列为
多中心、测试划分患者不重叠,但训练与测试中出现的中心相同。因此本研究不能
确立向未参与训练的机构或采集环境的泛化,那需要在训练中缺席的中心开展专门
的外部验证。

\textbf{测试队列构成与中心规模。}测试队列四个分期各含 60 名患者,而训练与
验证遵循队列分布。总体测试表现因此不应被解读为筛查人群的患病率加权估计。
四个测试中心贡献的患者数也极不均衡,从 2 到 115 不等;因此补充材料中的
逐中心估计仅是描述性的,最小中心的结果应谨慎解读。

\textbf{单一配置与单一种子。}比较基于一个预设种子,我们不将其作为运行间
变异性的证据;在主要终点上,联合配置与基线之间的差距相对单次运行所能分辨
的幅度而言很小,因此我们将联合配置表述为最佳观察结果,而非已确立的改进。
设计也仅在预设的 ResNet-50 框架下评估,其在其他骨干架构上的稳健性尚待确立。

\textbf{参考标准与病灶监督。}分级标签来自专家超声读片而非组织病理学,这对
此规模的无创筛查队列是必要的;因此模型校准于专家共识,而非组织级真值。
病灶分支由粗边界框监督;报告的 Dice 与 IoU 量化的是与这些弱标注的一致
程度,不是病灶分割精度的测量——本队列中不存在像素级参考标准。

\section{结论}
\label{sec:conclusion}

我们提出了 SynAP-Fib,一种面向血吸虫病相关肝纤维化的细粒度超声分级模型,
其中解剖部位与弱病灶证据经由门控、梯度隔离的残差注入,既充当训练时监督,
又充当推理时上下文,部署时无需解剖元数据。两个先验的 $2\times2$ 比较揭示
了一种非可加的实证模式:任一先验单独使用均使预设复合分级风险恶化,而其
组合则不然,并取得四种配置中最佳观察的复合风险。联合配置还以很小的计算
代价额外输出解剖部位与弱病灶定位,同时保持基线水平的分级表现。证明这些
收益能否在训练中缺席的种子、骨干与机构上成立,以及分离联合配置的收益究竟
来自辅助任务正则还是来自预测解剖上下文本身的注入,都仍是未来工作。
